\section{Descending properties of modules}
\label{section-descent}

\noindent
We address the faithfully flat descent of the properties from Theorem
\ref{theorem-projectivity-characterization} that characterize projectivity.
In the presence of flatness, the property of being a Mittag-Leffler module
descends:

\begin{lemma}
\label{lemma-ffdescent-ML}
\begin{reference}
Email from Juan Pablo Acosta Lopez dated 12/20/14.
\end{reference}
Let $R \to S$ be a faithfully flat ring map. Let $M$ be an $R$-module. If the
$S$-module $M \otimes_R S$ is Mittag-Leffler, then $M$ is Mittag-Leffler.
\end{lemma}

\begin{proof}
Write $M = \colim_{i\in I} M_i$ as a directed colimit
of finitely presented $R$-modules $M_i$.
Using Proposition \ref{proposition-ML-characterization}, we see that we
have to prove that for each $i \in I$ there exists $i \leq j$, $j\in I$
such that $M_i\rightarrow M_j$ dominates $M_i\rightarrow M$.

\medskip\noindent
Take $N$ the pushout
$$
\xymatrix{
M_i  \ar[r] \ar[d] & M_j \ar[d] \\
M \ar[r] & N
}
$$
Then the lemma is equivalent to the existence of $j$ such that
$M_j\rightarrow N$ is universally injective, see
Lemma \ref{lemma-domination-universally-injective}.
Observe that the tensorization by $S$
$$
\xymatrix{
M_i\otimes_R S  \ar[r] \ar[d] & M_j\otimes_R S \ar[d] \\
M\otimes_R S \ar[r] & N\otimes_R S
}
$$
Is a pushout diagram.  So because
$M \otimes_R S = \colim_{i\in I} M_i \otimes_R S$
expresses $M\otimes_R S$ as a colimit of $S$-modules of
finite presentation, and $M\otimes_R S$ is Mittag-Leffler,
there exists $j \geq i$ such that $M_j\otimes_R S\rightarrow N\otimes_R S$
is universally injective. So using that $R\rightarrow S$ is faithfully flat
we conclude that $M_j\rightarrow N$ is universally injective too.
\end{proof}

\begin{lemma}
\label{lemma-ffdescent-countable}
Let $R \to S$ be a faithfully flat ring map. Let $M$ be an $R$-module. If
the $S$-module $M \otimes_R S$ is countably generated, then $M$
is countably generated.
\end{lemma}

\begin{proof}
Say $M \otimes_R S$ is generated by the elements
$y_i$, $i = 1, 2, 3, \ldots$. Write $y_i = \sum_{j = 1, \ldots, n_i}
x_{ij} \otimes s_{ij}$ for some $n_i \geq 0$, $x_{ij} \in M$
and $s_{ij} \in S$. Denote $M' \subset M$ the submodule generated
by the countable collection of elements $x_{ij}$. Then
$M' \otimes_R S \to M \otimes_R S$ is surjective as the
image contains the generators $y_i$. Since $S$ is faithfully flat
over $R$ we conclude that $M' = M$ as desired.
\end{proof}

\noindent
At this point the faithfully flat descent of countably generated projective
modules follows easily.

\begin{lemma}
\label{lemma-ffdescent-countable-projectivity}
Let $R \to S$ be a faithfully flat ring map.  Let $M$ be an $R$-module.
 If the $S$-module $M \otimes_R S$ is countably generated and projective,
then $M$ is countably generated and projective.
\end{lemma}

\begin{proof}
Follows from Lemmas \ref{lemma-descend-properties-modules},
\ref{lemma-ffdescent-ML}, and \ref{lemma-ffdescent-countable} and
Theorem \ref{theorem-projectivity-characterization}.
\end{proof}

\noindent
All that remains is to use d\'evissage to reduce descent of projectivity in the
general case to the countably generated case.  First, two simple lemmas.

\begin{lemma}
\label{lemma-lift-countably-generated-submodule}
Let $R \to S$ be a ring map, let $M$ be an $R$-module, and let $Q$ be a
countably generated $S$-submodule of $M \otimes_R S$.  Then there exists a
countably generated $R$-submodule $P$ of $M$ such that
$\Im(P \otimes_R S \to M \otimes_R S)$ contains $Q$.
\end{lemma}

\begin{proof}
Let $y_1, y_2, \ldots$ be generators for $Q$ and write $y_j = \sum_k x_{jk}
\otimes s_{jk}$ for some $x_{jk} \in M$ and $s_{jk} \in S$. Then take $P$ be
the submodule of $M$ generated by the $x_{jk}$.
\end{proof}

\begin{lemma}
\label{lemma-adapted-submodule}
Let $R \to S$ be a ring map, and let $M$ be an $R$-module.  Suppose $M
\otimes_R S = \bigoplus_{i \in I} Q_i$ is a direct sum of countably generated
$S$-modules $Q_i$.  If $N$ is a countably generated submodule of $M$, then
there is a countably generated submodule $N'$ of $M$ such that $N' \supset N$
and $\Im(N' \otimes_R S \to M \otimes_R S) =
\bigoplus_{i \in I'} Q_i$ for some subset $I' \subset I$.
\end{lemma}

\begin{proof}
Let $N'_0 = N$.  We construct by induction an increasing sequence of countably
generated submodules $N'_{\ell} \subset M$ for $\ell = 0, 1, 2, \ldots$
such that: if $I'_{\ell}$ is the set of $i \in I$ such that the projection of
$\Im(N'_{\ell} \otimes_R S \to M \otimes_R S)$ onto $Q_i$ is
nonzero, then $\Im(N'_{\ell + 1} \otimes_R S \to M \otimes_R
S)$ contains $Q_i$ for all $i \in I'_{\ell}$.  To construct $N'_{\ell + 1}$
from $N'_\ell$, let $Q$ be the sum of (the countably many) $Q_i$ for
$i \in I'_{\ell}$, choose $P$ as in Lemma
\ref{lemma-lift-countably-generated-submodule}, and then let $N'_{\ell + 1} =
N'_{\ell} + P$.  Having constructed the $N'_{\ell}$, just take $N' =
\bigcup_{\ell} N'_{\ell}$ and $I' = \bigcup_{\ell} I'_{\ell}$.
\end{proof}

\begin{theorem}
\label{theorem-ffdescent-projectivity}
Let $R \to S$ be a faithfully flat ring map.  Let $M$ be an $R$-module.
 If the $S$-module $M \otimes_R S$ is projective, then $M$ is projective.
\end{theorem}

\begin{proof}
We are going to construct a Kaplansky d\'evissage of $M$ to show that it is a
direct sum of projective modules and hence projective.  By Theorem
\ref{theorem-projective-direct-sum} we can write $M \otimes_R S =
\bigoplus_{i \in I} Q_i$ as a direct sum of countably generated $S$-modules
$Q_i$.  Choose a well-ordering on $M$.  Using transfinite recursion we are going
to define an increasing family of submodules $M_{\alpha}$ of $M$, one for each
ordinal $\alpha$, such that $M_{\alpha} \otimes_R S$ is a direct sum of some
subset of the $Q_i$.

\medskip\noindent
For $\alpha = 0$ let $M_0 = 0$.  If $\alpha$ is a limit ordinal and $M_{\beta}$
has been defined for all $\beta < \alpha$, then define $M_\alpha =
\bigcup_{\beta < \alpha} M_{\beta}$.  Since each $M_{\beta} \otimes_R S$ for
$\beta < \alpha$ is a direct sum of a subset of the $Q_i$, the same will be
true of $M_{\alpha} \otimes_R S$.  If $\alpha + 1$ is a successor ordinal and
$M_{\alpha}$ has been defined, then define $M_{\alpha + 1}$ as follows.  If
$M_{\alpha} = M$, then let $M_{\alpha +1} = M$.  Otherwise choose the smallest
$x \in M$ (with respect to the fixed well-ordering) such that $x \notin
M_{\alpha}$. Since $S$ is flat over $R$, $(M/M_{\alpha}) \otimes_R S = M
\otimes_R S/M_{\alpha} \otimes_R S$, so since $M_{\alpha} \otimes_R S$ is
a direct sum of some $Q_i$, the same is true of $(M/M_{\alpha}) \otimes_R
S$.   By Lemma \ref{lemma-adapted-submodule}, we can find a countably generated
$R$-submodule $P$ of $M/M_{\alpha}$ containing the image of $x$ in
$M/M_{\alpha}$ and such that $P \otimes_R S$ (which equals $\Im(P
\otimes_R S \to M \otimes_R S)$ since $S$ is flat over $R$) is a
direct sum of some $Q_i$. Since $M \otimes_R S = \bigoplus_{i \in I} Q_i$
is projective and projectivity passes to direct summands, $P \otimes_R S$ is
also projective.  Thus by Lemma \ref{lemma-ffdescent-countable-projectivity},
$P$ is projective.  Finally we define $M_{\alpha + 1}$ to be the preimage of $P$
in $M$, so that $M_{\alpha + 1}/M_{\alpha} = P$ is countably generated and
projective.  In particular $M_{\alpha}$ is a direct summand of $M_{\alpha + 1}$
since projectivity of $M_{\alpha + 1}/M_{\alpha}$ implies the sequence $0
\to M_{\alpha} \to M_{\alpha + 1} \to
M_{\alpha + 1}/M_{\alpha} \to 0$ splits.

\medskip\noindent
Transfinite induction on $M$ (using the fact that we constructed
$M_{\alpha + 1}$ to contain the smallest $x \in M$ not contained in
$M_{\alpha}$) shows that
each $x \in M$ is contained in some $M_{\alpha}$.  Thus, there is some large
enough ordinal $S$ satisfying: for each $x \in M$ there is $\alpha \in S$ such
that $x \in M_{\alpha}$.  This means $(M_{\alpha})_{\alpha \in S}$ satisfies
property (1) of a Kaplansky d\'evissage of $M$.  The other properties are clear
by construction.  We conclude
$M = \bigoplus_{\alpha + 1 \in S} M_{\alpha + 1}/M_{\alpha}$.
Since each $M_{\alpha + 1}/M_{\alpha}$ is projective
by construction, $M$ is projective.
\end{proof}












